\documentclass[11pt,letterpaper]{article}
\usepackage[T1]{fontenc}
\usepackage{lmodern,amsmath,amssymb,amsthm,microtype}
\usepackage[margin=1in]{geometry}
\usepackage[hidelinks]{hyperref}
\providecommand{\tightlist}{\setlength{\itemsep}{0pt}\setlength{\parskip}{0pt}}
\DeclareUnicodeCharacter{2264}{\ensuremath{\le}}
\DeclareUnicodeCharacter{2265}{\ensuremath{\ge}}
\DeclareUnicodeCharacter{2260}{\ensuremath{\ne}}
\DeclareUnicodeCharacter{2212}{\ensuremath{-}}

\DeclareUnicodeCharacter{2211}{\ensuremath{\sum}}

\DeclareUnicodeCharacter{2208}{\ensuremath{\in}}
\title{Nine remaining full-leaf trace patterns in the four-isolate height-six Frankl case}
\author{Aafi Khan Malek\\ \small Independent researcher, Ahmedabad, India\\
\small\href{mailto:aafimalek2032@gmail.com}{aafimalek2032@gmail.com}}
\date{28 November 2026}
\begin{document}
\maketitle
This note continues the
\href{https://github.com/Aafimalek/novel-discoveries/tree/0e7028693b3320939007ed50f57f319520ca7c67/frankl-height-six/verification/exactly-four-isolates}{preceding four-isolate analysis}.
A weak family is finite, nonempty, union-closed, contains no empty member,
has no strict chain of seven members, and has every ground point in
at most half its members. Work on a smallest ground among such families.

\textbf{Written necessary conditions with an exact finite replay.} The
residual family is \emph{not} excluded. There is no new Lean claim.

Write p for the single source, L for the number of leaves, x for the low
isolate, and y,z,w for the three high isolates. Let B be the
source-avoiding members and b its size. The source-containing members
are the full star plus all 16 isolate traces except \{x\}. The previous
trace proof gives b odd, f\_B(x)=(b+1)/2, each high isolate occurs at
most (b−1)/2 times, and every leaf at most (b−15)/2 times. There are no
B-members with trace \{x\}. The full B-top is U without p: the source
has no distinct forcing predecessor. More explicitly, the checked
avoidance characterization notMem\_avoid\_iff\_forces says that a ground
point q is absent from avoid G r exactly when q forces r. Thus avoid G p
is U without p.~For every leaf l, its only distinct forcing predecessor
is p, so avoid G l is U without p and l: this is the checked lemma
avoid\_eq\_coPair\_of\_edge in FranklForcingFiber.lean, using weak
ground minimality. Each maximum belongs to G by avoid\_mem\_weak and
avoids p, so they are members of B. The B-top and the L co-leaf maxima
are all distinct and all have the full isolate trace.

\section{Co-pair forcing by filling a missing large
set}\label{co-pair-forcing-by-filling-a-missing-large-set}

The set C=U without p and x belongs to B. Otherwise add C to B. The
enlarged family is union-closed: every union of C with an old member is
C or the old full B-top. The new denominator is b+1. The x-frequency
does not change, and all other frequencies can rise by at most one, so
every point satisfies the weak half bound.

A chain containing C could be lifted to the original family by replacing
C with U without x and, if present above it, the full B-top with U. The
new chain is strict at every step since p is absent below C, and both
lifted sets are original members. Thus the enlargement does not violate
height six. Its ground still omits p, contradicting ground minimality
and proving C∈B.

\section{More general full-leaf trace
enlargement}\label{more-general-full-leaf-trace-enlargement}

Let M be the set of isolate traces of members of B containing all L
leaves. It is union-closed, contains the full isolate set and the high
triple by the preceding argument, and avoids \{x\}. Let V be the 15
trace sets other than \{x\}, and let W be any upper set in V ordered by
inclusion. Add to B all full-leaf sets having traces in A=W without M.
Closure holds: a union involving an added set still has all leaves, its
trace in V, and a trace above the added one. Hence it is either old or
added.

If there is a strict chain using an added set, leave the portion
strictly before its first added member unchanged. Lift that member and
all following members by adjoining p.~Each such member has all leaves
and trace in V, so its lift is an original source-containing member.
Strictness is preserved (p was absent before the first lift, and
adjoining it to full-leaf sets preserves their strict trace inclusions).
Thus the enlarged family has height at most six.

Let a=\textbar A\textbar, and f\_A be trace frequencies in the added
sets. The enlarged family would be a weak counterexample on smaller
ground if

\[
 2 f_A(x)+1\le a,\qquad 2 f_A(v)\le a+1\quad(v\in\{y,z,w\}).
\]

The inequalities use just the known B frequency margins. Each leaf is
present in all additions, but its original margin of 15 is sufficient,
since a≤15. The first inequality already implies a\textgreater0.
Therefore no candidate M may admit W satisfying these inequalities.

\section{Safe-trace refinement}\label{safe-trace-refinement}

Every trace T of B, even of a member with fewer than all leaves, must be
safe for M: T union R must belong to M for every R in M. This follows by
unioning the member with the full-leaf member for R. There is at least
one (in fact at least L+2) B-member containing x and exactly one high
isolate. Here is the counting argument, independently of the enlargement
cut.

Write \(A_j^+\) and \(A_j^-\) for the numbers of members of \(B\) with
\(x\) present and absent, respectively, and exactly \(j\) high isolates.
There are no \(A_0^+\) members. There is at most one \(A_0^-\) member,
since a two-chain of members with empty trace can be followed by a
maximal five-member chain in the source-containing traces: the full star
with no isolates, three successive high isolates, and then \(x\).
Counting the high-isolate incidences gives
\[
 A_1^+ \ge 2-A_0^-+A_3^++A_2^-+2A_3^-.
\]
Indeed, the \(x\)-present part contributes \((b+1)-A_1^++A_3^+\),
and the \(x\)-absent part contributes
\((b-1)/2-A_0^-+A_2^-+2A_3^-\). Their sum is at most \(3(b-1)/2\).
The \(B\)-top and every co-leaf maximum, omitting \(p\) and that leaf,
contribute to \(A_3^+\), so \(A_3^+\ge L+1\) and \(A_1^+\ge L+2\).
Consequently M must admit at least one safe trace of the
form \{x,y\}, \{x,z\}, or \{x,w\}.

\section{Weighted obstruction for the safe
traces}\label{weighted-obstruction-for-the-safe-traces}

Fix any nonempty subset J of the three high isolates and assign weight
one to those in J and zero to the others. Put h=\textbar J\textbar{} and
score each trace by the number of its elements in J. Let u be the
minimum score among safe traces containing x, and v the minimum among
nonempty safe traces avoiding x. These minima exist as the full isolate
trace and the x-free high triple are both in M and hence safe for M.

Write b=2a+1. The a+1 members of B containing x each have score at least
u, with L+1 known full-trace maxima each contributing the extra h−u. The
a members avoiding x include at most one with empty trace; all remaining
scores are at least v, and the forced co-pair C has score h,
contributing the extra h−v. Thus the weighted sum is at least

\[
 (u+v)a+u-v+(h-u)(L+1)+h-v.
\]

It is at most ha, since each high isolate occurs at most a times. If
u+v≥h, this is impossible. In fact L≥1 implies that the constant part is
at least 3h−u−2v≥0, with equality only at u=v=h. In that equality case
the coefficient u+v=2h already exceeds h, and a≥1 since the co-pair
exists. Therefore each of the seven nonempty choices J must satisfy
u+v\textless h.

\section{Exact replay and
limitation}\label{exact-replay-and-limitation}

Run python3 check\_full\_leaf\_traces.py. It tests all 15-bit M masks,
all upper sets, and the seven nonempty weight supports. There are 2,214
union-closed M patterns with both required tops; 567 survive
enlargement; 283 also admit a requisite safe x-and-one-high trace. The
weighted test leaves 28 labeled patterns, or nine up to permuting the
three high isolates. The log lists canonical masks for all nine. These
are trace patterns, not actually realizable family examples. The checks
do not finish the four-isolate single-source case.

Each trace is encoded by a bitmask on \((x,y,z,w)\), with weights \((1,2,4,8)\).
The remaining profiles, up to reordering the three high isolates, are:
\begin{center}
\small
\begin{tabular}{r l}
\textbf{No.}&\textbf{Traces in the profile}\\\hline
1&0, 2, 3, 4, 5, 6, 7, 8, 9, 10, 11, 12, 13, 14, 15\\
2&0, 2, 3, 4, 5, 6, 7, 8, 10, 11, 12, 13, 14, 15\\
3&0, 2, 3, 4, 6, 7, 8, 10, 11, 12, 13, 14, 15\\
4&2, 3, 4, 5, 6, 7, 8, 10, 11, 12, 13, 14, 15\\
5&2, 3, 4, 6, 7, 8, 10, 11, 12, 13, 14, 15\\
6&2, 3, 4, 6, 7, 10, 11, 12, 13, 14, 15\\
7&2, 4, 6, 7, 10, 11, 12, 13, 14, 15\\
8&2, 6, 7, 10, 11, 12, 13, 14, 15\\
9&2, 6, 7, 10, 11, 14, 15\\
\end{tabular}
\end{center}

\end{document}
